\documentclass{article}

% ── NeurIPS 2024 template ────────────────────────────────────────────────────
% Compile on Overleaf: upload this file alongside neurips_2024.sty
% (available in the NeurIPS 2024 LaTeX template package on neurips.cc)
\usepackage[final]{neurips_2024}

\usepackage[utf8]{inputenc}
\usepackage[T1]{fontenc}
\usepackage{hyperref}
\usepackage{url}
\usepackage{booktabs}
\usepackage{microtype}

\title{AI Frontier: An Interactive Dashboard for Mapping the LLM Landscape}

\author{
  Avo Sarkissian \\
  Electrical and Computer Engineering \\
  Northeastern University \\
  Boston, MA 02115 \\
  \texttt{avohsarkissian@gmail.com}
}

\begin{document}

\maketitle

\begin{abstract}
The number of commercially available large language models has grown to
the point where no single leaderboard covers all the dimensions that matter
for practical use: price, throughput, latency, context window size, and
benchmark performance. AI Frontier is an interactive dashboard that aggregates
live data for every hosted language model Artificial Analysis benchmarks —
155 of them across 31 providers at the time of writing — refreshed every hour,
and makes the full competitive landscape navigable through ten specialized
views. Counts move hourly as models are added and delisted; the dashboard's
header stats are authoritative. The system is written entirely in Python and
requires no JavaScript build step: the same chart code runs server-side under
Plotly Dash for local development and client-side under Pyodide (WebAssembly)
on the deployed static site. This report covers the data pipeline, composite
intelligence scoring methodology, visualization type selection, and the design
decisions that shaped the final tool.
\end{abstract}

\section{Introduction}

Selecting a large language model for a real application used to be
straightforward — a handful of options existed and the best one for most tasks
was obvious. That changed quickly. As of spring 2026, over 29 companies offer
models through hosted APIs, each making different tradeoffs in price,
capability, and latency. A model that scores 70 on reasoning benchmarks might
cost twenty times more per token than one that scores 65. A local model that
fits on a consumer GPU might outperform a paid API for a narrow task.
Without a way to see all of this together, model selection is largely
guesswork.

Existing leaderboards do one thing well. The LMSYS Chatbot Arena
\cite{lmsys} ranks models on human preference. The Open LLM Leaderboard
\cite{openllm} tracks open-weight models on academic benchmarks. But neither
shows live pricing, and most don't update on a daily or hourly schedule.
Provider dashboards show pricing but not quality scores. There is no single
tool that combines current benchmark performance, price, throughput, and
hardware requirements for both API-hosted and locally-runnable models in one
place.

AI Frontier fills that gap. The core idea is simple: pull current price and
performance data for every major model on a regular schedule and surface it
through enough visualization types that different kinds of users can answer
different kinds of questions. A developer wants to know the cheapest model
above a quality threshold. A researcher wants to see which provider is
consistently at the top. A student wants to know which models will fit in
their GPU's memory. Each of those questions requires a different view of the
same underlying data.

The project was built as part of EECE 5642 (Data Visualization) at
Northeastern University. The dashboard is publicly accessible and the full
source is open at
\url{https://github.com/Avo-Sarkissian/AI-Frontier-Database}.

\section{Data Sources and Pipeline}

\subsection{Cloud Model Data}

The primary data source is the Artificial Analysis API
\cite{artificialanalysis}, which tracks pricing and benchmark data for
commercially hosted LLMs. For each model it provides: price per million
tokens (input and output, at the cheapest available provider), median tokens
per second (throughput), median time-to-first-token (latency), context window
size, and a composite intelligence score. The API is queried once per hour by
a background thread that starts when the application launches.

Each successful API response is written to a CSV cache and also saved as a
timestamped snapshot in a history directory. Over 90 daily snapshots have
accumulated. No tab currently reads them: a price-over-time view was planned
and is not built, and the snapshots additionally span a change in the price
basis on 2026-05-06, so any such chart must not cross that boundary without
restating the earlier values.

A blended price per million tokens is computed from the raw input/output
prices using a 3:1 output-to-input ratio, approximating real-world usage
patterns. This weighting is ours, not Artificial Analysis's: AA publishes the
opposite blend (3 parts input to 1 part output), so prices here read roughly
1.9x theirs. Output-weighting is the defensible basis for agentic and
retrieval-augmented workloads, where generated tokens dominate spend. This collapses two numbers into one for display and sorting while
keeping the full per-direction figures available in the detail table.

\subsection{Local Model and Hardware Data}

The second data source pairs a live scrape of open-weight models with a
hand-maintained catalog of GPU hardware specifications. GPU memory bandwidth and VRAM capacity were
pulled from official product pages (NVIDIA, AMD, Apple, Intel). Model
parameter counts and quantization support were taken from HuggingFace model
cards \cite{huggingface}.

The GPU catalog is static — it gets updated when significant new hardware
generations are released — while the open-weight model list is scraped hourly
by data/local\_scraper.py on the same cadence as the hosted-model data. It backs the Run
Local tab, which estimates whether a given model fits in a selected GPU's
memory and how fast it would run.

\section{Intelligence Score and Benchmark Methodology}

The intelligence score shown throughout the dashboard is inherited from
Artificial Analysis, which aggregates results from multiple public
evaluations into a single 0--100 composite. The score covers four categories:

\begin{itemize}
  \item \textbf{Reasoning:} MMLU-Pro \cite{mmlu_pro} and GPQA \cite{gpqa},
        testing factual knowledge and multi-step deductive reasoning.
  \item \textbf{Coding:} LiveCodeBench \cite{livecodebench} and SciCode
        \cite{scicode}, evaluating writing and debugging real programs.
  \item \textbf{Math:} AIME and MATH-500 \cite{math500}, covering
        competition-style quantitative problems.
  \item \textbf{Knowledge:} Humanity's Last Exam \cite{hle}, a recently
        released benchmark of expert-level questions across scientific domains.
\end{itemize}

Using multiple benchmarks matters here. A model fine-tuned heavily on MMLU
training data might rank first on that test alone but fall to the middle on
LiveCodeBench. Averaging across categories rewards models that are
consistently capable rather than narrowly optimized, and it lets API-hosted
and local open-weight models be compared on the same scale.

The main limitation of composite scores is that they smooth over real
differences. Two models with the same score can have completely different
strength profiles — one excellent at coding and weak at math, the other the
reverse. The Compare tab (radar chart) exists specifically to surface these
within-score differences.

\section{Dashboard Design}

\subsection{Tab Structure}

The dashboard has ten tabs, each targeting a specific question:
\textbf{Overview} (cost vs. quality for every model),
\textbf{Agent Stack} (three-tier model recommendations for agentic workflows),
\textbf{Landscape} (treemap of the industry by provider),
\textbf{Rankings} (top 25 by intelligence, value, or speed),
\textbf{Compare} (radar chart for up to five models),
\textbf{Budget} (projected monthly API cost from a token volume),
\textbf{Table} (full sortable list of every tracked model),
\textbf{Run Local} (VRAM compatibility and speed estimation),
\textbf{Image Gen} (ELO-ranked image generation leaderboard),
and \textbf{Video Gen} (ranked video generation comparison, from a curated
dataset rather than a live scrape).

\subsection{Visualization Choices}

\textbf{Overview: bubble scatter with Pareto frontier.}
The main challenge with comparing 255 models simultaneously is that a ranked
list hides tradeoffs. The bubble scatter encodes three variables at once:
x-axis is log-scaled price, y-axis is quality score, and bubble size
represents throughput. Log scale on price is necessary because model prices
span roughly three orders of magnitude — from under \$0.01 to over \$15 per
million tokens — and a linear axis compresses most of the data into one
corner of the chart.

A dotted Pareto frontier connects models that are undominated in the
price-quality tradeoff: these are the points where any improvement in quality
requires paying more, and any cost reduction means accepting lower quality.
The frontier is computed in Pandas on each data load using a standard
non-dominated sorting pass over price and quality. Visually, points on or
near this curve represent the best practical choices for most real
applications.

Each provider is assigned both a color and a distinct marker shape. The
redundant encoding handles colorblind accessibility without requiring a
separate mode or a legend that distinguishes only by hue.

\textbf{Rankings: horizontal bars with tier separators.}
For one-dimensional ranked comparisons, horizontal bar charts are clearer
than scatter plots — the comparison axis is explicit and label legibility is
better horizontally. Tier separator lines are drawn wherever the score gap
between adjacent models exceeds a fixed threshold, marking meaningful
performance breaks rather than implying uniform spacing across the full range.

\textbf{Compare: radar chart.}
Head-to-head comparison across five dimensions (intelligence, speed,
affordability, context window, latency) is a natural use case for a radar
chart. Each axis is normalized so all five dimensions are comparable on the
same scale. The limitation of radar charts for this purpose is that
interpretation depends on axis ordering, which is fixed in the current
implementation. Reordering axes would change the visible shape of a model's
profile without changing any underlying values.

\textbf{Landscape: treemap.}
A treemap where tile area encodes model count per provider and color
intensity encodes average intelligence score gives a quick sense of industry
structure — which providers have the most models, which cluster at the top,
which are new entrants with a narrow catalog. The treemap is not suitable for
precise value comparison, but it communicates structural patterns that scatter
plots miss. A ranked bar chart directly below it allows precise comparison of
the same providers.

\textbf{Deferred: price-over-time line chart.}
The historical snapshot system was built to support this, and the market does
move fast enough to justify it. It is not shipped: the snapshots span a change
in the price basis on 2026-05-06, so a naive line would show a doubling that
never happened. Publishing a misleading trend is worse than publishing none.

\textbf{Run Local: VRAM compatibility calculator.}
This tab addresses a question that no existing public tool answers directly:
which open-weight models can I actually run on my hardware? The user selects a
GPU and a quantization level. The dashboard filters the catalog to models
whose estimated VRAM requirement fits the selected GPU's memory and displays
estimated inference speed in tokens per second.

VRAM requirements scale approximately with parameter count and bit-width.
The formula accounts for model weights, key-value cache overhead, and
activation memory. Quantization levels range from full precision (FP16/BF16)
through 8-bit and 4-bit integer approximations, each roughly halving or
quartering the memory footprint relative to the one above.

\subsection{Design Principles}

Three decisions carried through all ten tabs.

\textbf{Data-ink ratio.}
Gridlines, outer frames, and chart borders were removed wherever they add
visual weight without adding information \cite{tufte}. The dark background
with high-contrast data elements reduces perceived clutter compared to a
white background with equivalent data density — the background recedes and
the data elements stand out.

\textbf{Colorblind accessibility.}
The nine largest providers each get both a color and a distinct marker shape —
a redundant encoding that survives color-vision deficiency — kept consistent
across all ten tabs. The remaining providers share a neutral \enquote{Other}
series, because a categorical palette cannot separate thirty-plus categories
reliably for any viewer. A user who learns the color scheme (or the
shapes) can navigate between tabs without re-reading a legend each time.

\textbf{Appropriate scales.}
Price axes use log scale throughout. Context window sizes, which range from
8K to over 1 million tokens, also use log scale. Linear scales for heavily
right-skewed distributions compress the interesting variation at the low end,
where most models actually sit.

\subsection{Technical Stack}

The dashboard runs on Plotly Dash, which compiles Python component trees into
a live web interface without a JavaScript build step. All visualization logic
is in Python using Plotly's graph objects API; Dash handles HTTP routing,
callback wiring, and component state.

Choosing a single language for the full stack — data pipeline, computation,
and UI — simplified deployment. There is no Node.js dependency and no separate
API server. The project deployed on Render \cite{render} as a single Python
process until 2026-06-23, and now ships as a static site on GitHub Pages: the
same Python chart code is executed in the browser by Pyodide (WebAssembly),
and an hourly GitHub Actions job performs the scrape that browsers cannot
(the Artificial Analysis API sends no CORS headers).

The tradeoff moved with the architecture. Under Dash, interactions that a
React application would handle locally needed a server round-trip. Under
Pyodide there is no round-trip, but the browser downloads a 4 MB Python bundle
before the controls become live, so the site paints from pre-rendered figures
first and upgrades to interactive once that load completes.

\section{Contributions}

What AI Frontier adds relative to existing tools is the combination of live
data with multi-modal coverage. Static leaderboards go stale within days of a
new model release. Most pricing tools don't include benchmark scores. AI
Frontier scrapes both automatically, covers text, image, and video generation
in a single dashboard, and includes local open-weight models alongside API
services — a side-by-side comparison that does not exist elsewhere in a
unified interface.

The Agent Stack tab is a practical extension of the core data: given a
workflow type, it recommends a three-tier model setup (reasoning, balanced,
and fast tiers) and surfaces specific model options at each tier. For users
building agentic pipelines, this translates the raw leaderboard data into a
concrete configuration decision.

The full dataset pipeline and all visualization modules are open-source.
The dashboard has served live data continuously since March 2026.

\section{Future Work}

Several extensions are worth building. Per-category breakdown scores
(reasoning, coding, math, knowledge separately) are collected by Artificial
Analysis but not yet surfaced in the Compare tab — adding them would let users
make more targeted selections than a single composite number allows.

The Budget tab currently estimates cost from a token count without accounting
for prompt caching or multi-turn context accumulation. Both matter for
production API usage and are straightforward to model given the existing
pricing data.

Measured API latency by region is the other obvious addition. Published
time-to-first-token medians from provider spec sheets are not always
consistent with real-world performance, especially under load. The hourly
polling infrastructure is already in place; extending it to include synthetic
request timing would give a more complete picture of real operational cost.

\section*{References}

\begin{thebibliography}{10}

\bibitem{artificialanalysis}
Artificial Analysis.
\newblock Artificial Analysis API: LLM performance and pricing benchmarks.
\newblock \url{https://artificialanalysis.ai}, 2024.

\bibitem{lmsys}
W.-L. Chiang, L. Zheng, Y. Sheng, A.~N. Angelopoulos, T. Li, D. Li,
H. Zhang, B. Zhu, M. Jordan, J.~E. Gonzalez, and I. Stoica.
\newblock Chatbot arena: An open platform for evaluating LLMs by human
preference.
\newblock \textit{arXiv preprint arXiv:2403.04132}, 2024.

\bibitem{openllm}
Hugging Face.
\newblock Open LLM Leaderboard.
\newblock \url{https://huggingface.co/spaces/HuggingFaceH4/open_llm_leaderboard},
2023.

\bibitem{mmlu_pro}
Y. Wang, X. Ma, G. Zhang, Y. Ni, A. Chandra, S. Guo, W. Ren,
A. Arulraj, X. He, Z. Jiang, T. Li, M. Ku, K. Wang, A. Zhuang,
R. Fan, X. Yue, and W. Chen.
\newblock MMLU-Pro: A more robust and challenging multi-task language
understanding benchmark.
\newblock \textit{arXiv preprint arXiv:2406.01574}, 2024.

\bibitem{gpqa}
D. Rein, B.~L. Hou, A.~C. Stickland, J. Petty, R.~Y. Pang,
J. Dirani, J. Michael, and S.~R. Bowman.
\newblock GPQA: A graduate-level google-proof Q\&A benchmark.
\newblock \textit{arXiv preprint arXiv:2311.12022}, 2023.

\bibitem{livecodebench}
N. Jain, K. Han, A. Gu, W. Li, F. Yan, T. Zhang, S. Wang, A. Solar-Lezama,
K. Sen, and I. Stoica.
\newblock LiveCodeBench: Holistic and contamination free evaluation of large
language models for code.
\newblock \textit{arXiv preprint arXiv:2403.07974}, 2024.

\bibitem{scicode}
M. Tian, Y. Fang, J. Liu, Y. Yu, Y. Li, K. McAlpine, P. Heinzelman,
D. Lymperopoulos, J. Cheng, J. Smith, A. Shrivastava, Q. Xu, B. Zhou,
and H. Ji.
\newblock SciCode: A research coding benchmark curated by scientists.
\newblock \textit{arXiv preprint arXiv:2407.13168}, 2024.

\bibitem{math500}
D. Hendrycks, C. Burns, S. Kadavath, A. Arora, S. Basart, E. Tang,
D. Song, and J. Steinhardt.
\newblock Measuring mathematical problem solving with the MATH dataset.
\newblock In \textit{Proceedings of the Neural Information Processing Systems
Track on Datasets and Benchmarks}, 2021.

\bibitem{hle}
D. Phan, S. Pawar, D. Kambhampati, et al.
\newblock Humanity's Last Exam.
\newblock Scale AI, 2025.
\newblock \url{https://scale.com/hle}.

\bibitem{huggingface}
T. Wolf, L. Debut, V. Sanh, J. Chaumond, C. Delangue, A. Moi, P. Cistac,
T. Rault, R. Louf, M. Funtowicz, et al.
\newblock Transformers: State-of-the-art natural language processing.
\newblock In \textit{Proceedings of EMNLP: System Demonstrations},
pages 38--45, 2020.

\bibitem{tufte}
E.~R. Tufte.
\newblock \textit{The Visual Display of Quantitative Information}.
\newblock Graphics Press, 2nd edition, 2001.

\bibitem{render}
Render.
\newblock Render cloud application platform.
\newblock \url{https://render.com}, 2024.

\end{thebibliography}

\end{document}
